\begin{table}[t]
\centering
\caption{Bounded pilot metrics. Each row has a distinct unit and denominator.}
\label{tab:headline-metrics}
\begin{threeparttable}
\begin{tabularx}{\linewidth}{@{}L r l L@{}}
\toprule
Construct & Result & 95\% Wilson & Unit and denominator \\
\midrule
B2 Integration accessibility & 63.9\% & -- & 6 device--interface cases \\
B3 Deployment manifest & 54.2\% & -- & 3 deployment objects; 24 applicable cells \\
B4 Physical definitions & 67.0\% & -- & 4 resource definitions \\
B5 Observability & 52.5\% & -- & 4 execution/diagnostic cases \\
B6 Preflight preventability & 66.7\% & 20.8--93.9\% & 3 definite scenarios of 4 eligible \\
B7 Constraint discovery & 87.5\% & 52.9--97.8\% & 8 incomplete fields of 13 \\
B8 Fully aligned claims & 33.3\% & 9.7--70.0\% & 6 bounded claims \\
B9 Core context expansion & 2.0$\times$ & -- & 5 opening classes \\
B10 Actionable public outcome & 41.7\% & 19.3--68.0\% & 12 documentation cases \\
\bottomrule
\end{tabularx}
\begin{tablenotes}[flushleft]\footnotesize
\item Results demonstrate operationalization in selected cases. They do not estimate forum-wide or industry-wide rates.
\item Wilson intervals are descriptive uncertainty for the stated denominator, not population inference. A dash marks a row whose result is a mean of ordinal component scores or a ratio rather than a proportion, for which a binomial interval would be undefined.
\item Component means are taken over known cells only; unknown cells are excluded rather than scored as zero.
\end{tablenotes}
\end{threeparttable}
\end{table}
